\begin{center}
\begin{minipage}{0.98\linewidth}
\begin{tcolorbox}[colback=white,colframe=CStatement!25,boxrule=0.4pt,
  arc=1mm,left=3mm,right=3mm,top=3mm,bottom=3mm]
\centering
\textbf{先走到垂足，再沿同一方向走同样远}\par\medskip
\begin{tikzpicture}[font=\small,>=Stealth]

\begin{scope}[x=1.25cm,y=1.25cm]
\coordinate (P) at (0,3);\coordinate (H) at (1.44,1.08);\coordinate (Q) at (2.88,-.84);
\draw[CStatement,thick] (-.85,-.6375)--(3.8,2.85) node[above left] {$l:Ax+By+C=0$};
\draw[CExample,very thick,->] (P)--(H);
\draw[CTrap,very thick,->] (H)--(Q);
\draw[CProof] (1.296,1.272)--(1.488,1.416)--(1.632,1.224);
\draw[CExample,thick] (.624,1.968)--(.816,2.112);
\draw[CTrap,thick] (2.064,.048)--(2.256,.192);
\fill[CExample] (P) circle (2.2pt);\fill[CProof] (H) circle (2.2pt);\fill[CTrap] (Q) circle (2.2pt);
\node[above left] at (P) {$P(x_0,y_0)$};
\draw[gray] (H)--(2.05,1.15);\node[right] at (2.05,1.15) {$H$};
\node[below right] at (Q) {$P'(x',y')$};
\node[above right,text=CExample] at (.65,2.3) {$PH$};
\node[above right,text=CTrap] at (2.25,.3) {$HP'$};
\end{scope}
\end{tikzpicture}\par\medskip
$PP'\perp l$，$PH=HP'$；$H$ 同时是垂足和中点。一次位移到 $H$，两次相同位移到 $P'$。图按等比例的具体位置绘制。
\end{tcolorbox}
\end{minipage}
\end{center}
